% Version 0.28. Original paired counts with one remaining image-density input.
% Pinned source: 7cbf7a0c32fca1badf05931879a37c3f8a2b03b4.
\section{Joining the paired counts on the original arithmetic objects}
\label{sec:paired-count-installation}

\status{JX1 closes every join from the literal
\texttt{SymplecticFourImageDense} predicate to the actual
\texttt{CountII} and \texttt{CountIII}.  The catalogue and
rigidity are now supplied by CX and RX, while AX/DX/MX/PX
retain the same arithmetic datum and coefficient normalization.
Image-density itself remains an arithmetic input.
JX2 consequently reduces the selected height/window branch
to \texttt{CountI}, \texttt{CountV}, and that one predicate,
with the already stated prime-pair hypotheses retained.}

% BEGIN JX1
\begin{theorem}[JX1: one-input reconstruction of the actual paired counts]
\label{claim:JX1}
The reconstructed chains give the literal implication
\[
 \boxed{\quad
 \texttt{SymplecticFourImageDense}
      \ \Longrightarrow\
              \texttt{CountII}\wedge\texttt{CountIII}.
 \quad}                                                        \tag{JX.1}
\]
Precisely, for every primitive source character $\chi$ there
is a natural exponent $C_\chi$, chosen before
$h,t,\Omega,\ell,c,e,\psi$, for which both source cases have
finite strict Selmer length and
\[
 \frac{\ell_{c.O}V(\Omega,c,e,\psi)}{v_{c.O}(\ell)}
                  \log\ell
                     \le(\ell h.N e)^{C_\chi}.                 \tag{JX.2}
\]
Every host, character, local condition and coefficient
quantifier in \texttt{CountShape} is retained.
\end{theorem}
\begin{proof}
Fix one witness $h_A$ of the source image-density predicate.
The catalogue theorem CX9 supplies the actual
\texttt{NumberSp4Alt}, with constants chosen before the
varying prime, algebraic closure, complex embedding and
finite family.  RX6 supplies the actual
\texttt{RigiditySymplecticFour}, including its literal
ordinary and infinitesimal local conditions.  BX8
supplies the complete \texttt{CohomologyHyp}; RV7, BR7,
FP8 and SP1--SP6 supply the rank-one, residual,
presentation and weighted-intersection inputs in their
already proved scopes.  We verify how these suppliers
meet on one original datum, rather than infer the result
from their names.

Fix $\chi$ and then an original host $h$ with $h.\chi=\chi$,
a permissible $t$, an original $\Omega$, an odd $\ell$,
an actual coefficient ring $O=c.O$, and
$\psi\in h.X(c,e)$ in case II or III, where $0<e<t$.
For the moment suppose $O$ contains the roots of unity
required by both $n_{\mathrm{II}}$ and
$n_{\mathrm{III}}$.  AX1--AX5 construct an actual
$D:\texttt{Raw}\,\Omega\,c$ with \texttt{Good} and
\texttt{OfCase}, $D.e=e$ and the original character.
In particular $e\ge2$ for these cases.  Apply $h_A$
to exactly these choices.  It gives
\[
 D.\texttt{Fin0},\qquad D.\texttt{NoEnd},\qquad
 \exists\iota\ \forall r,\
    D.\texttt{EPt}(r)\wedge\operatorname{Irr}(r)
                  \Longrightarrow\texttt{IsAut4Alt}(\iota,r).
                                                               \tag{JX.3}
\]
Fix this one $\iota$ before using any irreducible point.
AX6--AX10 give the exact ordinary, level, dual and safe
admissibility properties of every such point for this
same $\iota$.  The auxiliary imaginary quadratic fields
used in RX and CX may vary as their proofs specify;
they do not change $\iota$, $D$, or the original
representations over $\mathbb Q$.

Let $n',p_1,m_0$ be the source's case parameters.
AX7 supplies their precise level and root-order data
\[
      L=\ell\,2^{n'}h.Np_1,\qquad
      k=(\ell-1)\,2^{n'}(2h.o),
                                                               \tag{JX.4}
\]
with $m_0\mid h.o$, the required \texttt{KillsMu} facts,
and the $k$-th power identities for
$D.\psi,D.\xi,D.\mathrm{om}$.  The alternatives for
$(n',p_1)$ are exactly those in (SP.48).
DX2, DX4 and DX8 supply the actual global/local
finiteness, framed dimensions and local base;
MX10 therefore constructs a \texttt{Model} over this
same $O$.  Apply RX7 to (JX.3), using the just
established level and case facts.  It gives
$D.\texttt{AllIrr}\,D.\texttt{RigidAt}$ on this $D$,
not only rigidity of a separate automorphic
representation.

The point, conjugacy, orbit and tangent comparisons
PX1--PX9 now supply all of (G2) in SP10.
The hypotheses \texttt{Fin0}, \texttt{NoEnd} and
actual irreducible rigidity used there have each been
supplied in (JX.3) and the preceding paragraph.
AX10 supplies all of (G1) under $h_A$.
CX8 gives exactly SP9's finite catalogue for the
original Galois level $L$.  Its use of an enlarged
auxiliary automorphic level is already included in
its fixed polynomial exponent; it changes neither
$L$ in (JX.4) nor a hypothesis about $D$.

For clarity, the numerical comparison reached by
these joins is the following.  Choose the source
prime $p=p_0(\ell h.Np_1)\ge5$ and arithmetic
Frobenius $\sigma$ at the original chosen place.
It is outside $D.S$ and prime to $L$.  The three
$k$-th roots obtained by evaluating the original
finite characters have Eisenstein trace
\[
 \operatorname{tr}D.r_0(\sigma)
   =\iota^{-1}(p^{e+1}z_1+p^ez_2+pz_3+1).
                                                               \tag{JX.5}
\]
CX9, or equivalently CX8 followed by SP9, gives
nonvanishing of every irreducible trace difference
and an integer $N_*\ne0$ which bounds the product
over every finite pairwise nonisomorphic family.
PX10 gives the actual $D.\texttt{Ineq}(\sigma)$.
Thus, with the original valuation and the exact
factor two,
\[
 \frac{\ell_O D.\mathrm{Sha}}{v_O(\ell)}\log\ell
 \le -2\log|N_*|_\ell
 \le 2\log|N_*|
 \le 2Ck[L(e+1)]^A\log p.                                     \tag{JX.6}
\]
The first inequality includes finiteness of the
left-hand module.  The original strict Selmer
module embeds in the one used for this bound:
at each added prime its unramified negative-twist
Frobenius scalar is a root of unity times $p^{-e}$,
hence differs from $1$, and multiplication by that
nonzero scalar difference on $E/O$ is surjective.
This is the precise SP10 \texttt{shaEq} inequality.
It does not assert equality of the two modules.

SP10's arithmetic absorption of (JX.6) uses only
the displayed $L,k,p$, the case formulas for $n'$,
and the host's already stated inequalities.
The new catalogue exponent $A$ is fixed in advance,
so that argument applies unchanged and gives one
$C_\chi$ in (JX.2) for both cases.  This is a choice
before any of the objects fixed in the second
paragraph, since $C_\chi$ depends only on $\chi$
and the fixed catalogue constants.

Finally remove the extra root assumption on $O$.
Adjoin enough $2$-power roots for
$2^{n_{\mathrm{II}}+n_{\mathrm{III}}}$ and take a
common finite coefficient overring $O'$ containing
the original $O$.  SP8 and RV1 prove all required
character transports: kernels, orders, conductors,
the twists $\chi_P$, parity, both case tests and
both integers $n_{\mathrm{II}},n_{\mathrm{III}}$
are unchanged.  They also prove equality of the
normalized Selmer length on the original $\Omega$.
Apply the preceding bound over $O'$ and return
using that equality.  The exponent does not
depend on $O'$ or its ramification.  The resulting
\texttt{CountShape} for the union of the two cases
restricts to each case, proving (JX.1) with the
exact original coefficient-ring quantifier.
\end{proof}
% END JX1

% BEGIN JX2
\begin{corollary}[JX2: the reduced height and window input]
\label{claim:JX2}
The selected cohomological branch has the sufficient input
\[
 \texttt{CountI}\wedge\texttt{CountV}
      \wedge\texttt{SymplecticFourImageDense}
                     \ \Longrightarrow\ \texttt{PointHV}.
                                                               \tag{JX.7}
\]
With the already stated prime-pair inputs of VW,
the same premise gives its original-field cohomological
window conclusion for every $C_t\ge1$.  The
$C_t=401$ fixed-pair consequence retains the cutoff
exponent $400$.
\end{corollary}
\begin{proof}
JX1 supplies \texttt{CountII} and \texttt{CountIII}
with the exact \texttt{CountShape} quantifiers.
Together with the two stated counts these are
precisely the four premises of HV10.
The auxiliary \texttt{CountVs}, \texttt{CountIV}
and \texttt{CountZ} have already been supplied
by HV5, HV6 and CZ6.  Applying HV10 gives the
literal \texttt{PointHV}.

The order of choices in that application is
unchanged.  For each $\chi$ and each prescribed
layer-depth function $j$, first choose the four
count exponents, then the finitely many
bad-prime layer constants and the prime
thresholds, and then the single $C$ in (HV.45).
Only afterwards vary the host, $t,B,E,\Omega$,
local models and finite prime sets.
Consequently its conclusion is the original
\[
  \texttt{VVplusLe}\bigl(E,t,(Nt)^C+CB\bigr)
                                                               \tag{JX.8}
\]
on every original CM field of every permissible
host, with all the same layer hypotheses.
No constant is chosen after observing the field
or finite prime set.

TW8 supplies the single constant record and the
simultaneous layer theorem, and OB8 supplies the
odd-minus bound with its constant $240$.
Substituting the newly obtained \texttt{PointHV}
in the already reconstructed VW3--VW4 proof gives
the claimed window and fixed-pair consequences,
with the same prime-pair inputs and rescaling.
This closes the propagation from (JX.7).
It does not assert that its three premises, or
the separate global family, arithmetic gap,
theta/local installation and metric/spectral
comparison, have been constructed.
\end{proof}
% END JX2

\paragraph{What has and has not been reduced.}
The original four-count statements remain valid
interfaces in all historical arguments.  The
current sufficient condition uses two actual
count suppliers and image-density because JX1
now supplies the other two counts from that
predicate alone.  IX traces a substantive part
of the image-density route and isolates its
remaining actual promodular realization and
automorphic seed inputs.  No logical equivalence
between that sufficient route and RH is claimed.
